\chapter{Épreuve chronométrée 8}\label{ch:iv1:timed-test-8}

\begin{ivtrains}
Il faut la traiter d'une seule traite, sans notes ni calculatrice.
\end{ivtrains}

\ivtestsection{A}{6}{40}

\begin{question}\label{iq:iv1:timed-test-8:1}
Lequel des deux est le plus probable~: obtenir au moins un six en 4 lancers d'un dé équilibré, ou au moins un double six en 24 lancers de deux dés équilibrés~? Donnez les deux probabilités. Vous pouvez utiliser $(\frac{35}{36})^{24} \approx 0.509$.
\end{question}

\begin{question}\label{iq:iv1:timed-test-8:2}
Un marcheur sur une grille carrée part de l'origine et effectue 4 pas, chacun vers le haut, le bas, la gauche ou la droite, avec une même probabilité. Avec quelle probabilité le marcheur revient-il à l'origine après les 4 pas~?
\end{question}

\begin{question}\label{iq:iv1:timed-test-8:3}
Une action verse un dividende de 4 l'année prochaine, et ce dividende croît à taux constant pour toujours. Si le rendement exigé est de 8~\% et que le cours vaut 100, quel taux de croissance le marché intègre-t-il~? Quel est le cours si le taux de croissance vaut 5~\%~?
\end{question}

\begin{question}\label{iq:iv1:timed-test-8:4}
Vous lancez trois dés équilibrés. Quelle est la probabilité que le plus grand et le plus petit des trois résultats diffèrent exactement de 5~?
\end{question}

\begin{question}\label{iq:iv1:timed-test-8:5}
$X$ et $Y$ sont indépendantes et uniformes sur $[0, 1]$. Quelle est la probabilité que $X + Y < 1$, sachant que $X < \frac12$~?
\end{question}

\begin{question}\label{iq:iv1:timed-test-8:6}
Un jeu de 28 dominos contient chaque paire non ordonnée des nombres de 0 à 6 exactement une fois, doubles compris. Vous en tirez deux au hasard. Quelle est la probabilité qu'ils aient un nombre en commun, de sorte qu'on puisse les placer bout à bout~?
\end{question}

\ivtestsection{B}{6}{60}

\begin{question}\label{iq:iv1:timed-test-8:7}
Vous tirez des nombres uniformes indépendants sur $[0, 1]$ et vous arrêtez lorsqu'un nombre est plus petit que le précédent. Quelle est l'espérance du nombre de tirages~?
\end{question}

\begin{question}\label{iq:iv1:timed-test-8:8}
Dans un match, A a besoin de 2 victoires supplémentaires et B de 3. Le match est interrompu maintenant, et l'enjeu est réparti proportionnellement à la probabilité de gagner de chaque joueur, les parties étant équitables. Quelle part revient à A~? Et si A gagne chaque partie avec la probabilité $0.6$~?
\end{question}

\begin{question}\label{iq:iv1:timed-test-8:9}
Pour tirer un point uniformément dans la boule unité en dimension 3, on tire dans le cube englobant et on rejette les points extérieurs à la boule. Quel est le taux d'acceptation, et quelle est l'espérance du nombre de tirages par point accepté~? Quel est le taux d'acceptation en dimension 10, sachant que le volume de la boule y vaut $\frac{\pi^5}{120}$~?
\end{question}

\begin{question}\label{iq:iv1:timed-test-8:10}
Un jeu de 52 cartes est mélangé. Quelle est l'espérance du nombre de rois qui apparaissent avant le premier as~?
\end{question}

\begin{question}\label{iq:iv1:timed-test-8:11}
Cinq personnes entrent dans un ascenseur au rez-de-chaussée d'un immeuble comportant 6 étages au-dessus. Chacune choisit un étage de façon indépendante et uniforme. Quelle est l'espérance du nombre d'arrêts~? Quelle est la probabilité que les cinq personnes se rendent à des étages différents~?
\end{question}

\begin{question}\label{iq:iv1:timed-test-8:12}
Deux régresseurs ont chacun une corrélation de $\frac12$ avec $y$, et une corrélation de $\frac12$ entre eux. Quel est le $R^2$ de la régression de $y$ sur les deux~? Quel est-il avec un seul régresseur~?
\end{question}
